\chapter{Ce que Jenga est, et ce qu'il n'est pas}

Jenga construit des projets C et C++. Vous décrivez ce que vous voulez obtenir ;
il détermine dans quel ordre compiler, avec quel compilateur, quels drapeaux, et
il produit l'exécutable, la bibliothèque, l'APK ou la page Web.

Ce chapitre dit ce qu'il fait, ce qu'il ne fait pas, et surtout \emph{par quel
mécanisme} --- parce que ce mécanisme explique presque tout le reste du livre.

\section{La phrase qui décide de tout}

\begin{nkretenir}
\textbf{Un fichier \texttt{.jenga} est du Python, et Jenga l'EXÉCUTE.}

Ce n'est pas une métaphore. Voici la ligne, dans le chargeur :

\begin{nkcode}{Jenga/Core/Loader.py:268}
exec(filePath.read_text(encoding='utf-8-sig'), globals_dict)
\end{nkcode}

Votre description de projet n'est pas analysée comme un format de données : elle
est \emph{lancée}. Les appels que vous écrivez --- \texttt{project()},
\texttt{files()}, \texttt{links()} --- sont de vraies fonctions Python qui
remplissent des objets en mémoire.
\end{nkretenir}

Trois conséquences immédiates, et vous les emploierez dès le chapitre 3 :

\begin{center}
\begin{tabular}{@{}ll@{}}
\toprule
\textbf{Vous pouvez} & \textbf{Parce que} \\
\midrule
mettre une boucle \texttt{for} dans un \texttt{.jenga} & c'est du Python \\
lire une variable d'environnement & c'est du Python \\
appeler une fonction que vous écrivez & c'est du Python \\
importer un module & c'est du Python \\
\bottomrule
\end{tabular}
\end{center}

\begin{nkpiege}
\textbf{Et la contrepartie, qu'il vaut mieux connaître tout de suite :} une
erreur dans votre \texttt{.jenga} est une erreur \emph{Python}, avec sa pile
d'appels. Une virgule oubliée ne donne pas « fichier de projet invalide » mais
\texttt{SyntaxError}.

\textbf{Un \texttt{.jenga} peut aussi faire n'importe quoi} --- effacer un
dossier, appeler le réseau. Il a tous les droits de l'interpréteur. Ne lancez
pas un \texttt{.jenga} que vous n'avez pas lu, exactement comme vous ne lanceriez
pas un script inconnu.
\end{nkpiege}

\section{Ce que Jenga fait}

\begin{center}
\begin{tabular}{@{}ll@{}}
\toprule
\textbf{Il s'occupe de} & \textbf{Concrètement} \\
\midrule
trouver les sources & \texttt{files(["src/**.cpp"])} \\
choisir le compilateur & selon la plateforme visée \\
ordonner les projets & d'après les dépendances déclarées \\
ne recompiler que le nécessaire & voir le chapitre 10 --- et son piège \\
produire l'artefact final & \texttt{.exe}, \texttt{.a}, \texttt{.apk}, \texttt{.wasm}\dots \\
empaqueter et signer & \texttt{jenga package}, \texttt{jenga sign} \\
lancer, tester, déboguer & \texttt{run}, \texttt{test}, \texttt{gdb} \\
configurer votre éditeur & \texttt{ide-setup}, \texttt{compile-flags} \\
\bottomrule
\end{tabular}
\end{center}

\begin{nkretenir}
\textbf{Douze systèmes cibles sont déclarés}, et le livre distingue soigneusement
ceux qui sont \emph{éprouvés} de ceux qui sont seulement \emph{déclarés} --- c'est
l'objet du chapitre 15.

\begin{nkcode}{Jenga/Core/Api.py:66-83}
class TargetOS(Enum):
    WINDOWS  LINUX  MACOS  ANDROID  IOS  TVOS  WATCHOS  IPADOS
    VISIONOS  WEB  PS4  PS5  XBOX_ONE  XBOX_SERIES  SWITCH
    HARMONYOS  FREEBSD  OPENBSD
\end{nkcode}
\end{nkretenir}

\section{Ce que Jenga n'est pas}

\begin{nkretenir}
\textbf{Ce n'est pas un gestionnaire de paquets.} Il ne télécharge pas vos
dépendances depuis un dépôt public. Une bibliothèque tierce est un dossier que
vous avez, et vous lui écrivez un \texttt{.jenga} --- ou vous la liez par
\texttt{links()}.

\textbf{Ce n'est pas un générateur de projets, et c'est la différence avec
CMake.} CMake produit un \texttt{Makefile} ou une solution Visual Studio, qui
construisent ensuite. Jenga \textbf{construit lui-même} : il appelle le
compilateur. (Il \emph{peut} générer un projet, par \texttt{jenga gen}, mais
c'est un service annexe, pas son fonctionnement.)

\textbf{Ce n'est pas un système de tâches générique.} Il ne remplace ni
\texttt{make} ni un script : il connaît la compilation, l'édition de liens et
l'empaquetage, et rien d'autre.
\end{nkretenir}

\begin{nkpiege}
\textbf{Il ne devine pas votre bibliothèque système.} Écrire \texttt{links(["X11"])}
suppose que \texttt{libX11} est installée et trouvable. Jenga passe le drapeau à
l'éditeur de liens ; c'est ce dernier qui échoue, et son message ne nomme pas
Jenga.

C'est la source la plus fréquente d'un « Jenga ne marche pas » qui n'est pas un
défaut de Jenga.
\end{nkpiege}

\section{Le vocabulaire, en cinq mots}

\begin{center}
\begin{tabular}{@{}ll@{}}
\toprule
\textbf{Mot} & \textbf{Ce qu'il désigne} \\
\midrule
\emph{workspace} & l'espace de travail : le fichier racine, et tout ce qu'il inclut \\
\emph{project} & une unité qui produit un artefact : un exe, une lib, un test \\
\emph{configuration} & \texttt{Debug}, \texttt{Release}\dots --- un jeu de réglages \\
\emph{toolchain} & une chaîne d'outils : compilateur, éditeur de liens, archiveur \\
\emph{filter} & une condition sous laquelle certains réglages s'appliquent \\
\bottomrule
\end{tabular}
\end{center}

\begin{nkretenir}
\textbf{Un workspace contient des projets ; un projet ne contient rien.} C'est la
seule hiérarchie. Il n'y a pas de sous-projet, pas de groupe : un projet dépend
d'un autre, et c'est ainsi qu'on structure.
\end{nkretenir}

\section{Ce que vous saurez à la fin de ce livre}

\begin{center}
\begin{tabular}{@{}rl@{}}
\toprule
\textbf{Partie} & \textbf{Ce qu'elle donne} \\
\midrule
I & installer Jenga, le trouver, construire un projet de A à Z \\
II & écrire un \texttt{.jenga} : sources, sorties, filtres, dépendances \\
III & comprendre ce que Jenga fait de votre description \\
IV & viser une plateforme --- et savoir laquelle est éprouvée \\
V & lancer, tester, empaqueter, signer, déployer \\
VI & lire un échec, éviter les pièges, écrire une chaîne d'outils \\
\bottomrule
\end{tabular}
\end{center}

\begin{nkretenir}
\textbf{Tout ce que ce livre affirme a été mesuré sur la version 2.4.0}, et les
sorties de terminal qu'il cite sont des \emph{captures réelles} --- pas des
exemples recopiés. Quand une chose n'a pas pu être vérifiée ici, il le dit.
\end{nkretenir}

\begin{nkbilan}
Jenga construit du C et du C++, et le fait lui-même plutôt que de générer un
autre système de construction. Sa décision fondatrice est qu'un fichier
\texttt{.jenga} est \textbf{du Python exécuté} : vous y disposez d'un langage
complet, et vous en héritez aussi les erreurs. Il n'est ni gestionnaire de
paquets, ni générateur de projets, ni ordonnanceur de tâches. Un workspace
contient des projets, et c'est toute la hiérarchie.
\end{nkbilan}

\begin{nkquestions}
\begin{enumerate}
\item Qu'est-ce qu'un fichier \texttt{.jenga}, techniquement ?
\item Citez deux choses que cela vous autorise, et une qu'il faut craindre.
\item En quoi Jenga diffère-t-il de CMake ?
\item Que se passe-t-il si vous écrivez \texttt{links(["X11"])} sans avoir la
      bibliothèque ?
\item Quelle est la seule hiérarchie de Jenga ?
\end{enumerate}
\end{nkquestions}

\begin{nkpratique}
Ouvrez un \texttt{.jenga} réel du dépôt --- par exemple
\nkfile{Applications/NkDames/NkDames.jenga} --- et relevez :

\begin{enumerate}
\item les constructions qui sont du \emph{Python} et non du vocabulaire Jenga
      (\texttt{import}, \texttt{def}, \texttt{if}, \texttt{os.getenv}) ;
\item le nombre de blocs \texttt{with filter(\dots)} ;
\item ce que le fichier ferait de différent si vous supprimiez la ligne
      \texttt{files([\dots])}.
\end{enumerate}
Répondez à la troisième \emph{avant} de l'essayer.
\end{nkpratique}

\begin{nkdemo}
Prouvez qu'un \texttt{.jenga} est bien exécuté, et non analysé.

Ajoutez, au tout début d'un \texttt{.jenga} de votre choix, la ligne :

\begin{nkcode}{}
print("JE SUIS EXECUTE, et il est", __import__("datetime").datetime.now())
\end{nkcode}

Lancez \texttt{jenga info}. La ligne doit apparaître, avec l'heure courante.

\textbf{Ce que la démonstration établit}, et c'est plus que « ça s'affiche » :
l'heure \emph{change} d'une exécution à l'autre. Un format de données ne peut pas
produire une valeur qui change ; seul du code exécuté le peut.

Retirez ensuite la ligne. \alerte{} Ne la laissez pas : un \texttt{print} dans un
\texttt{.jenga} s'exécute à \emph{chaque} commande, y compris celles dont vous
lisez la sortie avec un script.
\end{nkdemo}

\begin{nklogique}
Un collègue vous dit : « Jenga est cassé, mon projet ne compile pas, il me sort
une erreur Python ».

\begin{enumerate}
\item Quelles sont les deux familles d'erreur possibles, et laquelle cette
      description désigne-t-elle ?
\item Comment distingueriez-vous, en une commande, une erreur \emph{dans sa
      description} d'une erreur \emph{dans son code C++} ?
\item Il ajoute : « et ça marchait hier ». Qu'est-ce que cette information
      élimine, et qu'est-ce qu'elle n'élimine pas ?
\end{enumerate}
\end{nklogique}
